\documentclass[10pt]{article}

\usepackage[utf8]{inputenc}

\usepackage[T1]{fontenc}

\usepackage{amsmath}

\usepackage{amsfonts}

\usepackage{amssymb}

\usepackage[version=4]{mhchem}

\usepackage{stmaryrd}

\usepackage{mathrsfs}

\usepackage{bbold}



\begin{document}

точек равномерной сходимости ряда (1). Отсюда следует, что сумма всякого ряда непрерывных функций, сходящегося на нек-ром отрезке, непрерывна на всюду плотном множестве этого отрезка. Вместе с тем существует и сходящийся во всех точках отрезка ряд непрерывных функций такой, что точки, в к-рых он сходится не равномерно, образуют всюду плотное множество рассматриваемого отрезка.



Почлевное ивте г ри ров ание ра вноме рно сходящи их я рядов. Пусть $X=[a, b]$. Если члены ряда



\[

\sum a_{n}(x), x \in[a, b],

\]



интегрируемы по Риману (по Лебегу) на отрезке $[a, b]$, a ряд (4) равномерно сходится на этом отрезке, то его сумма $s(x)$ также интегрируема по Риману (по Лебегу) на $[a, b]$ в для любого $x \in[a, b]$ имеет место равенство



\[

\int_{a}^{x} s(t) d t=\int_{a}^{x}\left[\Sigma a_{n}(t)\right] d t=\Sigma \int_{a}^{x} a_{n}(t) d t,

\]



причем ряд, стоящий в правой части равенства, сходится равномерно на отрезке $[a, b]$.



В этой теореме нельзя заменить условие равномерной сходимости ряда (4) просто условием его сходимости на отрезке $[a, b]$, так как существуют сходящиеся на отрезке ряды даже непрерывных функций с непрерывной суммой, для к-рых неверна формула (5). Вместе с тем существуют различные ее обобщения. Ниже приведен результат для интеграла Стиліьеса.



Если $g(x)$ - возрастающая на отрезке $[a, b]$ функция, $a_{n}(x)$ - интегрируемы по Стилтьесу относительно функции $g(x)$, ряд (4) сходится равномерно на отрезке $[a, b]$, то сумма $s(x)$ ряда (4) также интегрируема по Стилтьесу относительно функции $g(x)$,



\[

\int_{a}^{x} s(t) d g(t)=\sum \int_{a}^{x} a_{n}(t) d g(t)

\]



и ряд, стоящий в правой части равенства, сходится равномерно на отрезке $[a, b]$.



Обобщается формула (5) и на функции многих переменных.



Условия почленного диффе ренци ров ания рядов в те раинах ра вном е р но й с хо дим ости. Если члены ряда (4) непрерывно дифференцируемы на отрезке $[a, b]$, ряд (4) сходится в нек-рой точке этого отрезка, а ряд, составленный из производных членов ряда (4), равномерно сходится на $[a, b]$, то сам ряд (4) также равномерно сходится на отрезке $[a, b]$, а его сумма $s(x)$ непрерывно дифференцируема на нем и



\[

\frac{d}{d x} s(x)=\frac{d}{d x} \sum a_{n}(x)=\sum \frac{d}{d x} a_{n}(x) .

\]



В этой теореме условие равномерной сходимости ряда, получающегося из данного почленным дифференцированием, нельзя заменить просто условием его сходимости на отрезке $[a, b]$, так как существуют равномерно сходяциеся на отрезке ряды непрерывно диффференцируемых функций, для к-рых ряды, получающиеся из них почленным дифференцированием, сходятся на отрезке, однако сумма исходного ряда либо недифференцируема на всем рассматриваемом отрезке, либо дифференцируема, но ее производная не равна сумме ряда из производных.



Таким образом, наличие свойства равномерной сходимости у рядов, так же как и свойства абсолютной сходимости (см. Абсолютно сходлщийся рлд), позволяет перенести на эти ряды нек-рые правила действия с конечными суммами: равномерная сходимость - возможность почленно переходить к пределу, почлевво интегриро- вать и дифференцировать ряды (см. формулы (3) (6)), а абсолютная сходимость - возможность переставлять члены ряда в любом порядке без изменения его суммы и перемножать ряды почленно.



Свойства абсолютной и равномерной сходимости функциональных рядов независимы друг от друга. Так, ряд



\[

\sum_{n=0}^{\infty} \frac{x^{2}}{\left(1+x^{2}\right)^{n}}

\]



поскольку все его члены не отрицательны, абсолютно сходится на всей числовой оси, но заведомо точка $x=0$ не является его точкой равномерной сходимости, т. к. его сумма



\[

s(x)=\left\{\begin{array}{c}

1+x^{2}, \text { если } x \neq 0, \\

0, \quad \text { если } x=0

\end{array}\right.

\]



разрывна в этой точке.



Ряд



\[

\sum_{n=1}^{\infty} \frac{(-1)^{n+1}}{x^{2}+n}

\]



равномерно сходится на всей действительности оси, но не сходится абсолютно ни в какой ее точке.



Лит см. при ст. Ряд.



РАВНОМЕРНОЕ ПРИБЛИЖЕНИЕ - то Же, что чебицевское приближение.



РАВНОМЕРНОЕ IРОСТРАНСТВО - множество с определенной на нем равномерной структурой. Р а вноме рная с т р укт у р а (р а в номе рн о с т ь) на множестве $X$ определяется заданием нек-рой системы $\mathfrak{U}$ подмножеств произведения $X \times X$. При этом система $\mathfrak{U}$ должна быть фильтром (т. е. для любых $V_{1}, V_{2} \in \mathfrak{A}$ пересечение $V_{1} \cap V_{2}$ также содержится в $\mathfrak{A}$, и если $W \supset V, V \in \mathfrak{A}$, то $W \in \mathfrak{A})$ и должна удовлетворять следующим аксиомам.



U1) Всякое множество $V \in \mathfrak{Y}$ содержит диагональ $\Delta=$ $=\{(x, x): x \in X\}$.



U2) Если $V \in \mathscr{A}$, то $V^{-1}=\{(y, x):(x, y) \in V\} \in \mathfrak{A}$.



$\mathrm{U} 3)$ Для любого $V \in \mathfrak{A}$ сушествует $W \in \mathfrak{A}$ такое, что $W \circ W \subset V$, где $W \circ W=\{(x, y):$ существует такое $є \in X$, что $(x, z) \in W$ и $(z, y) \in W\}$. Элементы $\mathfrak{A}$ наз. о к р у жен и я м и равномерности, определяемой системой $\mathfrak{A}$.



Равномерность на множестве $X$ может быть определена также путем задания на $X$ системы покрытий $(5$, удовлетворяющей следующим аксиомам.



C1) Если $\alpha \in(5$ и $\alpha$ вписано в покрытие $\beta$, то $\beta \in(5$.



С2) Для любых $\alpha_{1}, \alpha_{2} \in(5$ существует покрытие $\beta \in(5$, к-рое з в е 3 дн о в п и с а п о в $\alpha_{1}$ и в $\alpha_{2}$ (т. е. для любой точки $x \in X$ все элементы $\beta$, содержащие $x$, лежат в нек-рых әлементах $\alpha_{1}$ и $\alpha_{2}$ ). Покрытия, принадлежащие (5, наз. р а в н о е р ны м и п о к р ы т и ям и $X$ (относительно равномерности, определяемой системой (5).



Указанные два способа задания равномерной структуры әквивалентны. Напр., если равномерная структура на $X$ задана системой окружений $\mathfrak{U}$, то система $\sqrt{5}$ равномерных покрытий $X$ может быть построена так. Для всякого $V \in \mathbb{E}$ семейство $\alpha(V)=\{Y(x): x \in X\}$ (где $\widehat{V}(x)=\{y:(x, y) \in V\})$ является покрытием $X$. Покрытие $\alpha$ принадлежит (5) тогда и только тогда, когда существует вписанное в $\alpha$ покрытие вида $\alpha(V), V \in \mathfrak{A}$. Обратно, если (5 - система равномерных покрытий Р. п., систему окружений образуют множества вида $\bigcup\{H \times H$ : $: H \in \alpha\}, \alpha \in \mathcal{S}$, и всевозможные множества, их содержащие.



Равномерная структура на $X$ может быть задана также с помощью системы псевдометрик. Всякая равномерность на множестве $X$ порождает топологию: $T=\{G \subset X$ : для любой точки $x \in G$ существует такое $V \in \mathscr{A}$, что $V(x) \in G\}$



Свойства Р. п. являются обобщением равномерных свойств метрических пространств. Если $(X, \rho)-$ ме-





\end{document}